\chapter{Исправленные опечатки и замечания}

\section*{Опечатки оригинала, исправленные в тексте}
\begin{center}\small
\begin{tabular}{@{}r p{0.40\textwidth} p{0.40\textwidth}@{}}
\toprule
с. & в оригинале & исправлено\\ \midrule
3 & <<\dots в которых \texttt{\}}на находит приложение>> & <<\dots она находит\dots>>\\
7 & <<а в последствии>> & <<а впоследствии>>\\
11 & <<(обобщенном )>> & <<(обобщенном)>>\\
13 & <<прообраз множества $A$ при отображении $|$>> & <<\dots при отображении $f$>>\\
15 & Пример: $(A\cap\ov B\cap C)\cap(A\cap C)\equiv\dots$ & $(A\cap\ov B\cap C)\cup(A\cap C)\equiv\dots$ (иначе не проходит вынесение $A\cap C$ по $5^\circ$)\\
16 & в выкладке упр.~2: <<$([A\cap A)\cap(C\cap\ov C)]$>> & <<$[(A\cap A)\cap(C\cap\ov C)]$>>\\
18 & упр.~1 \S\,2: $A\cup[(B\setminus\ov A)\sd C]\equiv(A\cap B)\sd C$ (не тождество) & $A\cap[(B\setminus\ov A)\sd C]\equiv(A\cap B)\sd(A\cap C)$ (восстановлено; см. решение)\\
19 & доказательство (1.6): сомножитель $(B\cup C\cup C)$ на четвертом месте & $(B\cup C\cup B)$\\
21 & нет номера у упражнения <<Будем предполагать\dots>> & <<4.>>\\
24 & $f[\varphi_{\alpha_1},\dots,\varphi_{\alpha_n}]$ & $f[\varphi_{\alpha_1},\dots,\varphi_{\alpha_m}]$\\
24 & <<Заменим в формуле (1.12) функции $\varphi_i$ на $\psi_i$>> & <<Заменим в этой формуле, согласно (1.12), \dots>>\\
25 & определение композиционной замкнутости с оговоркой <<но имеются суперпозиции\dots, не принадлежащие $\M$>> & оговорка снята (иначе неверна фраза <<всякое замкнутое множество является композиционно замкнутым>>)\\
25 & <<множество $\M_a$ замкнуто>> & <<множество $T_a$ замкнуто>>\\
27 & $\varphi_0=\dfrac{x_1+x_2}{1+c^{-1}x_1x_2}$ & $\dfrac{x_1+x_2}{1+c^{-2}x_1x_2}$ (релятивистский закон сложения скоростей)\\
29 & упр.~4: <<при $H\cap\ov A\ne\emptyset$ $Q(A)\subset\M(H)$>> & <<\dots будет $Q(A)\subset\M(H)$>>\\
\bottomrule
\end{tabular}
\end{center}

Мелкие дефекты набора скана (сбитые индексы, разорванные черты
дополнения, точки вместо запятых) исправлены без отдельных отметок.

\section*{Замечания к исходным рабочим заметкам}
Решения упражнений \S\,3 основаны на заметках из каталога
\texttt{notes/claude/} (файлы \texttt{p21-*.md}). При переносе исправлено:
\begin{itemize}
\item В разминочной таблице заметки к упр.~2 формула
  $(A\cap B)\cup(\ov A\cap I)$ отмечена как самодвойственная (<<обе
  $\equiv\ov A\cup B$>>). Это неверно: $F\equiv\ov A\cup B$, а
  $F^*=(A\cup B)\cap\ov A\equiv\ov A\cap B$ (что, кстати, правильно
  вычислено ниже в той же заметке при проверке леммы (L)). В книге эта пара
  приведена лишь как пример \emph{двойственных} формул.
\item В решении упр.~4 добавлена прямая проверка самодвойственности
  формулы $F_n$ и замечание о вырождении при $n=2$; в упр.~6 "--- явная
  ссылка на конструкцию $F_n$; в остальном решения следуют заметкам.
\end{itemize}

\section*{Объем}
Набраны страницы 1--29 оригинала: с. 1--4 (титул, аннотация,
предисловие), с. 5--12 (введение), с. 13--29 (гл.~1, \S\S\,1--4, до
начала \S\,5 <<Машинный способ задания функций. Коммутативные диаграммы>>).
Страница 29 включена, чтобы завершить упражнения \S\,4.
